% !TeX program = lualatex
% !TeX TXS-program:compile = txs:///lualatex/[--shell-escape]
\documentclass{article}
\usepackage[paper=legal, ink=black, seed=944, background=pine, light=lamp, sloppiness=.92,wear=.9,angle=steep,framing=cropped]{cacography}
\usepackage{amsmath}
\begin{document}
\cacographypage{stains={coffee,smudge},props={pen,clip},edges=spiral}
\section{Alles vor der Klausur}
Morgen 08:30, Raum B-204. Tasche heute packen!\newline
Ausweis, Wasser, zwei Stifte, Taschenrechner, Ersatzbatterie.\newline
\cacographynote*{nicht verschlafen}
Mathe: Kettenregel, Produktregel, Integrationsgrenzen.\newline
Physik: Impulserhaltung, Energieerhaltung, Einheiten.\newline
Chemie: Stoffmenge, pH, Reaktionsgleichungen.\newline
Statistik: Standardfehler, Verteilung, Hypothesentest.\newline
CS: Laufzeit abschätzen, Schleifen von Hand durchgehen.\newline
Noch \crossout[line]{4} 3 Übungsaufgaben rechnen.\newline
\crossout[double]{Nachtschicht} Schlaf, wirklich.\newline
Wenn Panik: zuerst leichte Aufgaben, danach die schweren.\newline
\crossout[scribble]{alles lernen} nur Luecken schliessen.\newline
Stichworte, Formeln und Zwischenschritte sauber hinschreiben.\newline
Fragen bei Unklarheiten sofort markieren.\newline
Abends nicht noch zehn neue Themen anfangen.\newline
Busverbindung prüfen; 20 Min. Puffer.
\end{document}
